% ---------------------------------------------------------------------------
% CONCLUSION
% File: sections/06_conclusion.tex
% ---------------------------------------------------------------------------

\section{Conclusion}
\label{sec:conclusion}

\hspace{1cm}This study presented the first four-climate-zone comparative
optimization of integrated bifacial agrivoltaic and anaerobic digestion
systems, introducing a novel three-component extended Land Equivalent
Ratio as the optimization objective and a 4E analysis framework for
comprehensive performance assessment. The following conclusions are
supported by the validated simulation results.

\textbf{C1 --- The eLER framework confirms multi-output land
	productivity at all four climate zones.}
Optimised eLER values of 1.641 (Konya), 1.657 (Almer\'{i}a), 1.723
(Ouagadougou), and 1.764 (Freiburg) confirm that integrated APV--AD
systems use land more productively than three separate monocultures
at every site. The two-component $LER_{\text{2C}}$ ranges from 1.030
to 1.088 --- all sites exceed 1.0, confirming that the agrivoltaic
system alone outperforms two separate monocultures even without the
biogas contribution.

\textbf{C2 --- eLER is non-monotone with solar resource, driven by
	PAR saturation.}
Freiburg, with the lowest annual GHI (1\,150\,kWh\,m$^{-2}$),
achieves the highest eLER (1.764), exceeding Konya despite 43\% more
solar irradiance. The mechanism is PAR saturation: at high-irradiance
sites, open-field PAR frequently exceeds the tomato light saturation
point (\SI{174}{W\,m^{-2}}), inflating the $LER_{\text{crop}}$
denominator and penalising the yield ratio. This finding challenges
the conventional site-selection preference for maximum irradiance in
agrivoltaic deployment and reframes the optimisation problem as
matching solar resource to crop photosynthetic capacity rather than
simply maximising PV output.

\textbf{C3 --- Boundary convergence is a universal, climate-independent
	structural result.}
The optimizer converges to the minimum physically permissible row
spacing ($d_{\text{row}}^* = \SI{3.03}{m}$) and module height
($H_m^* = \SI{1.50}{m}$) at all four sites, despite GHI varying by
a factor of 1.8 and mean annual temperature varying by
17\,\si{\celsius}. This confirms that the eLER landscape has its
global maximum at the physical constraint boundary at all study
climates. Any mechanical access requirement beyond the minimum
physical clearance imposes an eLER penalty (approximately
0.04--0.06 per additional metre of row spacing in the earlier
version of the model; see Section~\ref{sec:discussion_boundary}
for a caveat on this figure) --- a design
trade-off that must be explicitly evaluated for each farming
system context.

\textbf{C4 --- Thermal self-sufficiency is confirmed at all sites
	with negligible heating allocation.}
ESR values of 453--2\,000 at the optimised designs confirm that
5\% of annual PV output is more than sufficient to cover all
digester heating demand at every site and in every season, including
January at Freiburg (winter T$_{\min} \approx -12$\,\si{\celsius}).
Zero biogas is consumed for self-heating, confirming that the full
biogas output is available for valorisation under all seven
economic scenarios.

\textbf{C5 --- The carbon credit scenario (S6) is the dominant
	economic strategy across all climate zones and market contexts.}
S6 raises IRR by 1.7--4.6 percentage points above the baseline
at all sites, achieving 12.7\% (Konya), 23.0\% (Almer\'{i}a),
29.9\% (Ouagadougou), and 17.6\% (Freiburg). Ouagadougou provides
the strongest investment case, with a payback period of 3.3 years
and CO$_2$ avoidance of 861\,tCO$_2$\,ha$^{-1}$\,yr$^{-1}$,
driven by the combination of high electricity price, optimal
Arrhenius BMP at 33.2\,\si{\celsius}, and high grid emission
factor. Extending carbon credit mechanisms to smallholder APV--AD
installations in Sub-Saharan Africa and Central Asia is identified
as the highest-leverage policy intervention for accelerating
deployment.

\textbf{C6 --- Growing-season ET denominator correction materially
	changes the water savings metric.}
Reporting $LER_{\text{water}}$ with the annual ET$_0$ denominator
--- as in prior literature --- underestimates the agronomic water
benefit by a factor of 1.3$\times$ (Konya, 189-day season) to
3.5$\times$ (Ouagadougou, 122-day season). The corrected
growing-season values of 15.6--22.4\% represent 81--159\,mm of
irrigation water saved per growing season --- a quantity equivalent
to extending the growing season by 2--4 weeks or irrigating an
additional 16--22\% of land area. This correction is particularly
consequential for Sub-Saharan Africa, where prior annual-denominator
values of $\sim$4\% do not justify irrigation infrastructure
investment, while the corrected seasonal value of 15.6\% does.

\textbf{C7 --- All results are conservative lower bounds.}
Both model validations exhibit systematic conservative bias: the crop
yield model (V1 MAPE = 2.4\%) has a near-zero mean bias, while the
BMP temperature model (V2 MAPE = 26.9\% in the operating range)
systematically underestimates methane yield at temperatures below
30\,\si{\celsius}. All reported eLER values, biogas yields, NPV,
IRR, and CO$_2$ avoidance figures are therefore lower bounds on
true system performance, with cold-climate sites (Konya, Freiburg)
having the largest upside potential from improved BMP modelling.

\bigskip
\noindent The integrated APV--AD framework developed in this study
provides a replicable, climate-adaptive design tool for food-energy-water
nexus planning across the semi-arid-to-temperate spectrum. The open-source
simulation code, optimised design parameters, and full results database
are available as supplementary material to facilitate replication and
extension to additional crop species, climate zones, and policy contexts.